% Generated by jobs/build_results_table.py -- do not hand-edit.
% Table 2: ablation ladder. Row 1 is the frozen baseline of Table 1, same numbers.
\begin{table}[t]
\centering
\small
\setlength{\tabcolsep}{4pt}
\begin{tabular}{cccccccc}
\toprule
\multicolumn{3}{c}{Adapted} & \multicolumn{2}{c}{Reconstruction} & \multicolumn{3}{c}{Temporal (texel-space)} \\
\cmidrule(lr){1-3}\cmidrule(lr){4-5}\cmidrule(lr){6-8}
Cross-attn & Self-attn & Temporal & PSNR\,$\uparrow$ & SSIM\,$\uparrow$ & Flicker\,$\downarrow$ & Accel.\,$\downarrow$ & Drift \\
\midrule
\xmark & \xmark & \xmark & 11.50 & 0.3363 & 0.02074 & 0.03102 & 0.3621 \\
\cmark & \xmark & \xmark & 23.47 & 0.7209 & 0.00854 & 0.01175 & 0.2288 \\
\cmark & \cmark & \xmark & 24.50 & 0.7536 & 0.00839 & 0.01148 & 0.2213 \\
\cmark & \cmark & \cmark & \textbf{24.55} & \textbf{0.7554} & \textbf{0.00707} & \textbf{0.00771} & 0.2221 \\
\bottomrule
\end{tabular}
\caption{\textbf{Each component buys a different thing.} Same 24 objects and
same frozen row as Table~\ref{tab:comparison}. Adapting cross-attention alone recovers
most of both quantities (+11.97\,dB,
58.8\% flicker). Extending the adapter to self-attention adds
+1.03\,dB, better on 24/24 objects
($p<10^{-6}$), but its effect on coherence is not significant --- lower flicker on only
15/24 objects, $p=0.31$: it is a fidelity component, not a temporal
one. The temporal component is the reverse, lowering flicker 15.7\%
on 24/24 objects and acceleration 32.8\% on
24/24 (both $p<10^{-6}$) while PSNR moves
+0.05\,dB. Drift is a guard and is never marked best.
Cross-attn is rung~19; self-attn adds rungs 27--30 (28--30 add KL regularisers to
rung~27); temporal is rungs 31--34 together with every arm using the temporal blend.}
\label{tab:ablation}
\end{table}
